% Generated by convert_myst_to_latex.py from chapter08/chapter08_sec05.md.
% Change the master and convert again instead of editing this file.
\RequirePackage[l2tabu, orthodox, abort]{nag}
\documentclass[paper=a4, fontsize=11pt, parskip=half]{scrartcl}

% exercises_preamble.tex
%
% Präambel für die Übungsblätter zum Selbststudium (Kapitel N.3 bzw. N.5) als
% PDF, übernommen aus der Präambel der CA-2-Unterlagen
% (internal/exams/exam_2026ws/ca02/preamble.tex). Die Studierenden üben damit
% das Format von CA 3: Aufgaben auf Papier, gelöst lokal am PC.
% Die .tex-Dateien erzeugt convert_myst_to_latex.py aus dem Master.

% ── Seitenlayout ─────────────────────────────────────────────────────────
\usepackage[a4paper, width=150mm, height=240mm]{geometry}

\usepackage{amsmath, amssymb}

\usepackage[no-math]{fontspec}
\usepackage{unicode-math}
\usepackage{libertinus}

\usepackage{setspace}
\usepackage{microtype}
\usepackage{textcomp}
\usepackage[main=english]{babel}

% ✩ für den Schwierigkeitsgrad (Libertinus hat das Zeichen nicht)
\usepackage{pifont}
\usepackage{newunicodechar}
\newunicodechar{✩}{\ding{73}}

% ── Kopf-/Fußzeilen ──────────────────────────────────────────────────────
\usepackage[automark]{scrlayer-scrpage}
\clearpairofpagestyles

\ifoot{S. Gramsch}
\cfoot{TH Mannheim -- Applied Numerical Mathematics -- WS 2026/27}
\ofoot{\pagemark}

% Keine Titelei: Die erzeugte .tex-Datei beginnt mit der Überschrift
% "N.M Exercises" als \section*, darunter "How to Practice" und "Exercises"
% als \subsection*.

% ── Listen ───────────────────────────────────────────────────────────────
\usepackage{tabularx}
\usepackage{booktabs}
\usepackage{enumitem}

% ── Farben ───────────────────────────────────────────────────────────────
\usepackage{xcolor}
\definecolor{itwm_blue_04}{HTML}{005A94}
\definecolor{frankfurtBlueLight}{HTML}{AECAE2}

\definecolor{my_darkgray}{HTML}{484949}
\definecolor{my_lightgray}{HTML}{F3F4F4}
\definecolor{my_orange}{HTML}{E87846}

% Farben der Admonition "tip" im Buch (MyST book-theme, heller Modus):
% Rand und Symbol --myst-color-success, Kopfzeile --myst-color-success-bg
\definecolor{tip_green}{HTML}{16A34A}
\definecolor{tip_bg}{HTML}{F0FDF4}

% ── Python-Code ──────────────────────────────────────────────────────────
\usepackage{listings}

% Zeilenumbrüche im Code mit Pfeil markieren, damit beim Abtippen klar ist,
% dass es eine einzige Python-Zeile ist
\lstdefinestyle{python}{
  language=Python,
  basicstyle=\ttfamily\small\setstretch{1},
  keywordstyle=\color{itwm_blue_04}\bfseries,
  stringstyle=\color{my_orange!80!black},
  commentstyle=\color{gray},
  showstringspaces=false,
  columns=fullflexible,
  keepspaces=true,
  upquote=true,
  breaklines=true,
  postbreak=\mbox{\textcolor{gray}{$\hookrightarrow$}\space},
  extendedchars=true,
}

\lstdefinestyle{ausgabe}{
  basicstyle=\ttfamily\footnotesize\setstretch{1},
  columns=fullflexible,
  keepspaces=true,
  upquote=true,
}

\lstdefinestyle{inline}{
  basicstyle=\ttfamily,
  showstringspaces=false,
  columns=fullflexible,
  keepspaces=true,
  upquote=true,
}

% Code im Fließtext: \py|derivative(values, dt)|
\newcommand{\py}{\lstinline[style=inline]}

% ── Kästen ───────────────────────────────────────────────────────────────
\usepackage[most]{tcolorbox}
\tcbuselibrary{listings}

%% Aufgabe ---------------------------------------------------------------
% \begin{aufgabe}{Exercise 7.1 (✩)}{approx.\ 7 min}: Bezeichnung links,
% Richtzeit rechts (in CA 2 und CA 3 stehen dort die Punkte).
% Aussehen wie die Admonition "tip" im Buch: grüner Rand links, hellgrüne
% Kopfzeile mit grünem Stift, weißer Inhalt, leichter Schatten.
% Eine Aufgabe wird nicht über zwei Seiten geteilt: Passt sie nicht mehr,
% beginnt sie auf der nächsten Seite. Nur für Aufgaben, die länger als eine
% Seite sind, setzt der Konverter die teilbare Variante aufgabelang ein.
\tcbset{
  aufgabenstil/.style 2 args={
    enhanced,
    boxrule=0pt,
    frame hidden,
    borderline west={1.5pt}{0pt}{tip_green},
    title={\textcolor{tip_green}{\ding{46}}\enspace #1 \hfill #2},
    coltitle=black,
    fonttitle=\bfseries,
    colbacktitle=tip_bg,
    colback=white,
    drop fuzzy shadow=black!25,
    arc=1mm,
    left=4mm, right=4mm,
    before skip=4mm, after skip=4mm,
  },
}
\newtcolorbox{aufgabe}[2]{aufgabenstil={#1}{#2}}
\newtcolorbox{aufgabelang}[2]{aufgabenstil={#1}{#2}, breakable}

%% Code und Ausgabe ------------------------------------------------------
% breakable: längerer vorgegebener Code darf über eine Seite hinausgehen,
% sonst bleibt vor ihm eine halbe Seite leer
\newtcblisting{code}{
    listing only,
    enhanced,
    breakable,
    boxrule=0pt,
    frame hidden,
    sharp corners,
    colback=my_lightgray,
    left=4pt, right=4pt, top=2pt, bottom=2pt,
    listing options={style=python},
}

\newtcblisting{ausgabe}{
    listing only,
    enhanced,
    boxrule=0pt,
    frame hidden,
    sharp corners,
    colback=white,
    borderline west={2pt}{0pt}{my_lightgray!70!black},
    left=6pt, right=4pt, top=0pt, bottom=0pt,
    listing options={style=ausgabe},
}

% ── Verweise ─────────────────────────────────────────────────────────────
\usepackage[hidelinks]{hyperref}


\begin{document}
\section*{8.5 Exercises}

\subsection*{How to Practice}

\begin{itemize}[noitemsep]
  \item Open a new Jupyter notebook on your own computer, switch off the
        internet connection and do not use AI tools.
  \item Type the given code yourself and pay attention to indentation,
        parentheses and quotation marks. A gray arrow
        \textcolor{gray}{$\hookrightarrow$} marks a line that is too long for
        the page. Type it as one line.
  \item Write your predictions and explanations in Markdown cells.
\end{itemize}

\subsection*{Exercises}

These exercises are meant for self-study at home and review the material of Chapters 8.1 to 8.4. Allow about 90 minutes for them. Before starting the exercises, run the following cell. It contains the functions \py|central()| and \py|gradient_descent()| from Chapter 8.1 as well as \py|gradient()| and \py|gradient_descent_2d()| from Chapter 8.3. Several exercises use the same names, for example \py|cost|. If you go back to an earlier exercise, run its cells again in order. Also keep in mind that the methods stop when \py|max_iterations| is reached, and that a small slope alone does not prove that you have found a minimum.

The difficulty level is given in the title of each exercise:

\begin{itemize}
  \item ✩ Comprehension: predict and explain code and outputs (approx. 5 min)
  \item ✩✩ Application: write your own code and interpret results (approx. 10 min)
  \item ✩✩✩ Mini-project: combine several concepts from this chapter (approx. 30 min)
\end{itemize}

These times only cover working on the exercises. Allow extra time for reading, checking your results and debugging.

\begin{code}
import numpy as np
import matplotlib.pyplot as plt

# functions from Chapter 8.1
def central(function, point, dx):
    """Central difference quotient of function at the given point."""
    return (function(point + dx) - function(point - dx)) / (2 * dx)

def gradient_descent(function, start, learning_rate, tolerance,
                     max_iterations):
    """Search for a minimum of function, starting at start.

    Returns all visited points as an array.
    """
    point = start
    history = [point]
    for k in range(max_iterations):
        slope = central(function, point, 1e-6)
        # stopping criterion: slope almost zero
        if abs(slope) < tolerance:
            break
        point = point - learning_rate * slope
        history.append(point)
    return np.array(history)

# functions from Chapter 8.3
def gradient(function, point, dx):
    """Gradient of function at point = [x, y].

    Each partial derivative with the central difference quotient.
    """
    shift_x = np.array([dx, 0.0])
    shift_y = np.array([0.0, dx])
    derivative_x = (function(point + shift_x)
                    - function(point - shift_x)) / (2 * dx)
    derivative_y = (function(point + shift_y)
                    - function(point - shift_y)) / (2 * dx)
    return np.array([derivative_x, derivative_y])

def gradient_descent_2d(function, start, learning_rate, tolerance,
                        max_iterations):
    """Search for a minimum of function, starting at start = [x, y].

    Returns all visited points as an array, row = [x, y].
    """
    point = start
    history = [point]
    for k in range(max_iterations):
        slope = gradient(function, point, 1e-6)
        # stopping criterion: length of the gradient almost zero
        if np.linalg.norm(slope) < tolerance:
            break
        point = point - learning_rate * slope
        history.append(point)
    return np.array(history)
\end{code}

\clearpage
\begin{aufgabe}{Exercise 8.1 (✩)}{approx.\ 5 min}
In milling, the cost per part depends on the \textbf{feed rate}, i.e. on how fast the tool moves through the material. If the feed rate is low, machining takes a long time. If it is high, the tool wears out quickly. In a simplified model, the cost in € per part is $C(v) = (v - 4)^2 + 10$ with the feed rate $v$ in m/min. The model is only valid near $4\,\text{m/min}$. The starting value $v = 0$ in the code just illustrates how the method works. The following code is given:

\begin{code}
def cost(feed):
    """Cost in € per part for the feed rate in m/min."""
    return (feed - 4)**2 + 10

# three steps of gradient descent for three learning rates
for learning_rate in [0.25, 0.5, 1.0]:
    feed = 0.0
    history = [feed]
    for k in range(3):
        slope = central(cost, feed, 1e-6)
        feed = feed - learning_rate * slope
        history.append(feed)
    print(learning_rate, np.round(history, 2))
\end{code}

Write down your guess in a Markdown cell before you run the code.

\begin{enumerate}
  \item At which feed rate is the cost lowest, and how high is it there?
  \item The derivative is $C'(v) = 2 \cdot (v - 4)$. How large is the slope at the start $v = 0$?
  \item Which three feed rates does the code compute for each of the three learning rates? What do you notice about the behavior with a learning rate of $1.0$?
  \item Run the code and check your predictions.
\end{enumerate}

\end{aufgabe}

\clearpage
\begin{aufgabe}{Exercise 8.2 (✩)}{approx.\ 5 min}
A classmate retyped the gradient descent code for the order quantity from Chapter 8.1. Here is their code:

\begin{code}
def total_cost(quantity):
    """Annual cost C(Q) in € for the order quantity Q = quantity."""
    return 12000 / quantity * 50 + quantity / 2 * 1.2

quantity = 500
for k in range(3):
    slope = central(total_cost, quantity, 1e-6)
    quantity = quantity + 200 * slope
    cost = total_cost(quantity)
    print(f'Order quantity {quantity:8.1f} units, cost €{cost:.1f}')
\end{code}

Write down your guess in a Markdown cell before you run the code.

\begin{enumerate}
  \item Find the error.
  \item The slope at 500 units is $-1.8$ €/unit. Where does the first step land? Does the cost go up or down?
  \item After the second step, the order quantity is negative, and so is the cost. Why does \py|total_cost()| return a value at all, and why is it meaningless?
  \item Run the code and check your predictions.
\end{enumerate}

\end{aufgabe}

\clearpage
\begin{aufgabe}{Exercise 8.3 (✩)}{approx.\ 5 min}
A \textbf{snap-action switch} has two stable positions. The stored energy in mJ depends on the deflection $x$ in mm: $E(x) = x^4 - 4x^2 + x$. The curve has two valleys, at about $x = -1.47\,\text{mm}$ and $x = 1.35\,\text{mm}$, with a hump at about $x = 0.13\,\text{mm}$ in between. When released, the switch settles in whichever valley it rolls into. The following code is given:

\begin{code}
def energy(x):
    """Stored energy in mJ at the deflection x in mm."""
    return x**4 - 4 * x**2 + x

for start in [0.5, -0.5, 0.1]:
    history = gradient_descent(energy, start, 0.02, 1e-4, 500)
    x_end = history[-1]
    e_end = energy(x_end)
    print(f'Start {start:5.2f} mm: end at {x_end:6.3f} mm, {e_end:6.3f} mJ')
\end{code}

Write down your guess in a Markdown cell before you run the code.

\begin{enumerate}
  \item In which valley does the method end up for each of the three starting values? Pay particular attention to the starting value $0.1\,\text{mm}$.
  \item Each valley is a \textbf{local minimum}. The deepest valley is also called the \textbf{global minimum}. Which valley is the global minimum? Use the formula to estimate.
  \item How can we find the global minimum of a function with several valleys if gradient descent is the only tool we have?
  \item Run the code and check your predictions.
\end{enumerate}

\end{aufgabe}

\clearpage
\begin{aufgabe}{Exercise 8.4 (✩)}{approx.\ 5 min}
A spring is deflected by $1$, $2$ and $3\,\text{cm}$, and the measured forces are $2.3$, $4.1$ and $6.2\,\text{N}$. The model is $F = c \cdot s$ with the spring stiffness $c$ in N/cm. The following code is given:

\begin{code}
deflection = np.array([1.0, 2.0, 3.0])
force = np.array([2.3, 4.1, 6.2])

for c in [2.0, 2.1]:
    deviation = force - c * deflection
    mse = np.mean(deviation**2)
    print(c, np.round(deviation, 2), round(mse, 4))
\end{code}

Write down your guess in a Markdown cell before you run the code.

\begin{enumerate}
  \item For both values of $c$, compute the three deviations and the mean squared error by hand.
  \item Which spring stiffness fits better? What is the unit of the mean squared error?
  \item For $c = 2.0$, all deviations are positive. What does this mean for the model line compared with the data points?
  \item Run the code and check your predictions.
\end{enumerate}

\end{aufgabe}

\clearpage
\begin{aufgabe}{Exercise 8.5 (✩✩)}{approx.\ 10 min}
A truck drives $400\,\text{km}$ on the autobahn. The faster it drives, the shorter the trip and the lower the driver's wages. Fuel consumption, however, increases with speed. At a speed of $v$ in km/h, the truck uses $20 + 0.0015 \cdot v^2$ liters per $100\,\text{km}$. The driver costs €30 per hour, and a liter of diesel costs €1.70. So the trip costs

\[
C(v) = \frac{400}{v} \cdot 30 + \frac{400}{100} \cdot (20 + 0.0015 \cdot v^2) \cdot 1.70.
\]

\begin{enumerate}
  \item Write a function \py|trip_cost(v)| and plot the cost for speeds from $50$ to $100\,\text{km/h}$.
  \item Determine the cheapest speed with \py|gradient_descent()|. Start at $60\,\text{km/h}$ with a learning rate of $10$, a tolerance of $10^{-3}$ and at most 500 iterations. What happens with a learning rate of $50$?
  \item In Germany, trucks over $7.5\,\text{t}$ may drive at most $80\,\text{km/h}$ on the autobahn. How much more does a trip at $80\,\text{km/h}$ cost than one at the cheapest speed?
\end{enumerate}

\end{aufgabe}

\clearpage
\begin{aufgabe}{Exercise 8.6 (✩✩)}{approx.\ 10 min}
In Chapter 8.3, we claimed that the weighted center of gravity minimizes the sum of the squared distances. Let us check this with gradient descent. The plants and loads from Chapter 8.3 are:

\begin{code}
plants = np.array([[6.0, 24.0], [16.0, 30.0], [20.0, 12.0],
                   [8.0, 6.0], [44.0, 10.0]])
loads = np.array([12, 10, 10, 8, 6])
\end{code}

\begin{enumerate}
  \item Write a function \py|effort_squared(location)| that computes $\sum n_i \cdot d_i^2$.
  \item Start \py|gradient_descent_2d()| at $(40, 30)$ with the learning rate of $0.2$ from Chapter 8.3, a tolerance of $0.01$ and at most 200 iterations. What happens? Then try the learning rates $0.02$, $0.01$ and $0.005$.
  \item Compare the result with the weighted center of gravity. Why does this cost function need a much smaller learning rate than \py|effort()| in Chapter 8.3?
\end{enumerate}

\end{aufgabe}

\clearpage
\begin{aufgabe}{Exercise 8.7 (✩✩✩) Mini-project: Precise robot positioning}{approx.\ 30 min}
In Chapter 5.4, we determined the joint angles of the SCARA robot for the placement site $P_1 = (0.30, 0.35)$ with a grid search in $1°$ steps. After evaluating $27\,331$ positions, the gripper with $\alpha = 9°$ and $\beta = 99°$ was still $3.2\,\text{mm}$ away from the site. Now we will do better. For this exercise, you only need the following formula from Chapter 5.4 for the position of the gripper with the arm lengths $l_1 = 0.4\,\text{m}$ and $l_2 = 0.3\,\text{m}$:

\[
x = l_1 \cos\alpha + l_2 \cos(\alpha + \beta), \qquad
y = l_1 \sin\alpha + l_2 \sin(\alpha + \beta).
\]

Joint 2 may only bend between $0°$ and $150°$.

\textbf{Part 1:} Write a function \py|gripper(angles)| that returns the position of the gripper as an array for \py|angles = [alpha, beta]| in degrees. \py|np.deg2rad()| converts degrees to radians. As a check, in the home position $(0°, 0°)$, the gripper should be at $(0.7, 0)$. Then write a cost function \py|cost(angles)| that computes the square of the distance between the gripper and $P_1$ in $\text{mm}^2$. Use it to check the result of the grid search.

\textbf{Part 2:} Start \py|gradient_descent_2d()| at the result of the grid search with a learning rate of $0.01$, a tolerance of $10^{-3}$ and at most 500 iterations. Which angles do you find, and how far is the gripper from $P_1$ now? How many times was \py|cost()| called, compared with the grid search?

\textbf{Part 3:} Now start the method in the position $(90°, 0°)$, in which the arm points straight up, and store the result in \py|angles_up|. Which angles does it find? The following code draws both solutions as a line from the origin through joint 2 to the gripper. It expects the solution from Part 2 in \py|angles_best|, the target $P_1$ in \py|target| and the length of arm 1 in \py|l_1|.

\begin{code}
# both positions of the robot arm
fig, ax = plt.subplots(figsize=(6, 6))
for angles, name in [(angles_best, 'solution from Part 2'),
                     (angles_up, 'solution from Part 3')]:
    alpha = np.deg2rad(angles[0])
    joint_2 = l_1 * np.array([np.cos(alpha), np.sin(alpha)])
    tip = gripper(angles)
    ax.plot([0, joint_2[0], tip[0]], [0, joint_2[1], tip[1]],
            marker='o', linewidth=3, label=name)
ax.scatter(target[0], target[1], marker='x', s=100, color='black',
           label='P1')
ax.set_aspect('equal')
ax.set_xlabel('x in m')
ax.set_ylabel('y in m')
ax.set_title('Two positions for the same placement site')
ax.legend()
ax.grid(True)
plt.show()
\end{code}

\textbf{Closing question:} Is the robot allowed to move to the solution from Part 3? What does this mean for using gradient descent, and how can grid search and gradient descent be combined effectively?

\end{aufgabe}

\end{document}
